\subsection{Perturbation of Fredholm maps}

THEOREM 1. --- Let E and F be Banach spaces. The set \(\mathcal{F}(\mathrm{E};\mathrm{F})\) of Fredholm operators from E to F is open in the Banach space \(\mathcal{L}(\mathrm{E};\mathrm{F})\) , and the map \(u\mapsto\mathrm{ind}(u)\) of \(\mathcal{F}(\mathrm{E};\mathrm{F})\) into \(\mathbf{Z}\) is locally constant.

The theorem follows from the following more precise statement:

PROPOSITION 4. --- Let E and F be Banach spaces, \(u_{0}\colon\mathrm{E}\to\mathrm{F}\) a Fredholm operator and \(v_{0}\) a quasi-inverse of \(u_{0}\) . There exists an open neighborhood U of \(u_{0}\) in \(\mathcal{L}(\mathrm{E};\mathrm{F})\) and an analytic mapping \(\varphi\colon\mathrm{U}\to\mathcal{L}(\mathrm{F};\mathrm{E})\) such that

\begin{enumerate}
\def\labelenumi{(\roman{enumi})}
\item
  One has \(\varphi(u_{0}) = v_{0}\) ;
\item
  For every \(u\) in \(U\) , the map \(\varphi(u)\) is a quasi-inverse of \(u\) , and in particular \(u\) is a Fredholm map;
\item
  For every \(u\) in U, we have \(\operatorname{ind}(u) = \operatorname{ind}(u_0)\).
\end{enumerate}

By prop. 2 of III, p. 42, (ii), there exist topological direct sum decompositions \(\mathrm{E} = \mathrm{E}_1 \oplus \mathrm{E}_2\) and \(\mathrm{F} = \mathrm{F}_1 \oplus \mathrm{F}_2\) , and there exists \(\alpha_0 \in \mathcal{I}(\mathrm{E}_1; \mathrm{F}_1)\) such that \(\mathrm{E}_2\) and \(\mathrm{F}_2\) are finite-dimensional and \(u_0\) is represented by the matrix \(\left( \begin{array}{cc}\alpha_0 & 0\\ 0 & 0 \end{array} \right)\) relative to these decompositions.

Let U be the set of elements \(u\) of \(\mathcal{L}(\mathrm{E};\mathrm{F})\) such that, in the matrix representation \(\left( \begin{array}{cc}\alpha & \beta \\ \gamma & \delta \end{array} \right)\) of \(u\) with respect to these decompositions, we have \(\alpha \in \mathcal{I}(\mathrm{E}_1;\mathrm{F}_1)\) . Since \(\mathcal{I}(\mathrm{E}_1;\mathrm{F}_1)\) is open in \(\mathcal{L}(\mathrm{E}_1;\mathrm{F}_1)\)

(prop. 3 of III, p. 57), U is an open neighborhood of \(u_0\) in \(\mathcal{L}(\mathrm{E};\mathrm{F})\) . For \(u = \begin{pmatrix} \alpha & \beta \\ \gamma & \delta \end{pmatrix}\) in U, set

\[
\varphi (u) = v _ {0} + \left( \begin{array}{c c} \alpha^ {- 1} - \alpha_ {0} ^ {- 1} & 0 \\ 0 & 0 \end{array} \right).\tag{1}
\]

One has \(\varphi(u_{0}) = v_{0}\) and the map \(\varphi\) is analytic (loc. cit.). Modulo continuous linear maps of finite rank, we have the congruences

\[
u \equiv \left( \begin{array}{c c} \alpha & 0 \\ 0 & 0 \end{array} \right) , v _ {0} \equiv \left( \begin{array}{c c} \alpha_ {0} ^ {- 1} & 0 \\ 0 & 0 \end{array} \right) , \varphi (u) \equiv \left( \begin{array}{c c} \alpha^ {- 1} & 0 \\ 0 & 0 \end{array} \right).
\]

Therefore, \(\varphi(u)\) is a quasi-inverse of \(u\) . Every element \(u\) from U defines by restriction an isomorphism from \(E _1\) onto a topological complement of \(F _2\) in F. By prop. 3 of III, p. 44, we therefore have

\[
\begin{array}{c} \operatorname{ind} (u) = \operatorname{codim} _ {\mathrm{F}} (u (\mathrm{E} _ {1})) - \operatorname{codim} _ {\mathrm{E}} (\mathrm{E} _ {1}) \\ = \dim (\mathrm{F} _ {2}) - \dim (\mathrm{E} _ {2}) = \operatorname{ind} (u _ {0}), \end{array}
\]

which concludes the proof.
% semantic-fragments: legacy-input-seam
\relax{}
